% Options for packages loaded elsewhere
\PassOptionsToPackage{unicode}{hyperref}
\PassOptionsToPackage{hyphens}{url}
\PassOptionsToPackage{dvipsnames,svgnames,x11names}{xcolor}
\documentclass[
  11pt,
]{article}
\usepackage{xcolor}
\usepackage[a4paper,margin=2.2cm]{geometry}
\usepackage{amsmath,amssymb}
\setcounter{secnumdepth}{-\maxdimen} % remove section numbering
\usepackage{iftex}
\ifPDFTeX
  \usepackage[T1]{fontenc}
  \usepackage[utf8]{inputenc}
  \usepackage{textcomp} % provide euro and other symbols
\else % if luatex or xetex
  \usepackage{unicode-math} % this also loads fontspec
  \defaultfontfeatures{Scale=MatchLowercase}
  \defaultfontfeatures[\rmfamily]{Ligatures=TeX,Scale=1}
\fi
\usepackage{lmodern}
\ifPDFTeX\else
  % xetex/luatex font selection
  \setmainfont[]{Palatino}
  \setmonofont[Scale=0.82]{Menlo}
\fi
% Use upquote if available, for straight quotes in verbatim environments
\IfFileExists{upquote.sty}{\usepackage{upquote}}{}
\IfFileExists{microtype.sty}{% use microtype if available
  \usepackage[]{microtype}
  \UseMicrotypeSet[protrusion]{basicmath} % disable protrusion for tt fonts
}{}
\makeatletter
\@ifundefined{KOMAClassName}{% if non-KOMA class
  \IfFileExists{parskip.sty}{%
    \usepackage{parskip}
  }{% else
    \setlength{\parindent}{0pt}
    \setlength{\parskip}{6pt plus 2pt minus 1pt}}
}{% if KOMA class
  \KOMAoptions{parskip=half}}
\makeatother
\usepackage{color}
\usepackage{fancyvrb}
\newcommand{\VerbBar}{|}
\newcommand{\VERB}{\Verb[commandchars=\\\{\}]}
\DefineVerbatimEnvironment{Highlighting}{Verbatim}{commandchars=\\\{\}}
% Add ',fontsize=\small' for more characters per line
\newenvironment{Shaded}{}{}
\newcommand{\AlertTok}[1]{\textcolor[rgb]{1.00,0.00,0.00}{\textbf{#1}}}
\newcommand{\AnnotationTok}[1]{\textcolor[rgb]{0.38,0.63,0.69}{\textbf{\textit{#1}}}}
\newcommand{\AttributeTok}[1]{\textcolor[rgb]{0.49,0.56,0.16}{#1}}
\newcommand{\BaseNTok}[1]{\textcolor[rgb]{0.25,0.63,0.44}{#1}}
\newcommand{\BuiltInTok}[1]{\textcolor[rgb]{0.00,0.50,0.00}{#1}}
\newcommand{\CharTok}[1]{\textcolor[rgb]{0.25,0.44,0.63}{#1}}
\newcommand{\CommentTok}[1]{\textcolor[rgb]{0.38,0.63,0.69}{\textit{#1}}}
\newcommand{\CommentVarTok}[1]{\textcolor[rgb]{0.38,0.63,0.69}{\textbf{\textit{#1}}}}
\newcommand{\ConstantTok}[1]{\textcolor[rgb]{0.53,0.00,0.00}{#1}}
\newcommand{\ControlFlowTok}[1]{\textcolor[rgb]{0.00,0.44,0.13}{\textbf{#1}}}
\newcommand{\DataTypeTok}[1]{\textcolor[rgb]{0.56,0.13,0.00}{#1}}
\newcommand{\DecValTok}[1]{\textcolor[rgb]{0.25,0.63,0.44}{#1}}
\newcommand{\DocumentationTok}[1]{\textcolor[rgb]{0.73,0.13,0.13}{\textit{#1}}}
\newcommand{\ErrorTok}[1]{\textcolor[rgb]{1.00,0.00,0.00}{\textbf{#1}}}
\newcommand{\ExtensionTok}[1]{#1}
\newcommand{\FloatTok}[1]{\textcolor[rgb]{0.25,0.63,0.44}{#1}}
\newcommand{\FunctionTok}[1]{\textcolor[rgb]{0.02,0.16,0.49}{#1}}
\newcommand{\ImportTok}[1]{\textcolor[rgb]{0.00,0.50,0.00}{\textbf{#1}}}
\newcommand{\InformationTok}[1]{\textcolor[rgb]{0.38,0.63,0.69}{\textbf{\textit{#1}}}}
\newcommand{\KeywordTok}[1]{\textcolor[rgb]{0.00,0.44,0.13}{\textbf{#1}}}
\newcommand{\NormalTok}[1]{#1}
\newcommand{\OperatorTok}[1]{\textcolor[rgb]{0.40,0.40,0.40}{#1}}
\newcommand{\OtherTok}[1]{\textcolor[rgb]{0.00,0.44,0.13}{#1}}
\newcommand{\PreprocessorTok}[1]{\textcolor[rgb]{0.74,0.48,0.00}{#1}}
\newcommand{\RegionMarkerTok}[1]{#1}
\newcommand{\SpecialCharTok}[1]{\textcolor[rgb]{0.25,0.44,0.63}{#1}}
\newcommand{\SpecialStringTok}[1]{\textcolor[rgb]{0.73,0.40,0.53}{#1}}
\newcommand{\StringTok}[1]{\textcolor[rgb]{0.25,0.44,0.63}{#1}}
\newcommand{\VariableTok}[1]{\textcolor[rgb]{0.10,0.09,0.49}{#1}}
\newcommand{\VerbatimStringTok}[1]{\textcolor[rgb]{0.25,0.44,0.63}{#1}}
\newcommand{\WarningTok}[1]{\textcolor[rgb]{0.38,0.63,0.69}{\textbf{\textit{#1}}}}
\usepackage{longtable,booktabs,array}
\usepackage{caption}
\captionsetup[table]{skip=6pt}
\newcounter{none} % for unnumbered tables
\usepackage{calc} % for calculating minipage widths
% Correct order of tables after \paragraph or \subparagraph
\usepackage{etoolbox}
\makeatletter
\patchcmd\longtable{\par}{\if@noskipsec\mbox{}\fi\par}{}{}
\makeatother
% Allow footnotes in longtable head/foot
\IfFileExists{footnotehyper.sty}{\usepackage{footnotehyper}}{\usepackage{footnote}}
\makesavenoteenv{longtable}
\setlength{\emergencystretch}{3em} % prevent overfull lines
\providecommand{\tightlist}{%
  \setlength{\itemsep}{0pt}\setlength{\parskip}{0pt}}
\usepackage{booktabs}\usepackage{longtable}\usepackage{fvextra}\DefineVerbatimEnvironment{Highlighting}{Verbatim}{breaklines,commandchars=\\\{\}}\newfontfamily\symfont{Arial Unicode MS}\usepackage{newunicodechar}\newunicodechar{→}{{\symfont →}}\newunicodechar{⇔}{{\symfont ⇔}}\newunicodechar{✓}{{\symfont ✓}}\newunicodechar{✗}{{\symfont ✗}}
\usepackage{bookmark}
\IfFileExists{xurl.sty}{\usepackage{xurl}}{} % add URL line breaks if available
\urlstyle{same}
% fallback for those not using the hyperref driver hyperxmp:
\makeatletter
\@ifundefined{xmpquote}{\newcommand{\xmpquote}[1]{#1}}{}
\makeatother
\hypersetup{
  colorlinks=true,
  linkcolor={blue},
  filecolor={Maroon},
  citecolor={Blue},
  urlcolor={Blue},
  pdfcreator={LaTeX via pandoc}}

\author{}
\date{}

\begin{document}

\section{Real-Time Operating Systems --- Final
Examination}\label{real-time-operating-systems-final-examination}

\textbf{KMUTNB · Faculty of Engineering · M.Eng. in Electrical \&
Computer Engineering}

{\def\LTcaptype{none} % do not increment counter
\begin{longtable}[]{@{}
  >{\raggedright\arraybackslash}p{(\linewidth - 2\tabcolsep) * \real{0.5000}}
  >{\raggedright\arraybackslash}p{(\linewidth - 2\tabcolsep) * \real{0.5000}}@{}}
\toprule\noalign{}
\endhead
\bottomrule\noalign{}
\endlastfoot
\textbf{Course} & 010135311 --- Real-Time Operating Systems \\
\textbf{Semester} & 1/2026 \\
\textbf{Coverage} & Weeks 1--5 (foundations, task model, RMS, EDF,
aperiodic/sporadic servers) and Week 6 up to and including mutexes (race
conditions, critical sections, queues, semaphores, mutexes) \\
\textbf{Duration} & 3 hours (180 minutes) \\
\textbf{Format} & \textbf{Open book} \\
\textbf{Total} & 100 points \\
\end{longtable}
}

Name: \_\_\_\_\_\_\_\_\_\_\_\_\_\_\_\_\_\_\_\_\_\_\_\_\_\_\_\_\_\_
Student ID: \_\_\_\_\_\_\_\_\_\_\_\_\_\_\_\_\_\_\_\_ Date:
\_\_\_\_\_\_\_\_\_\_\_\_

\begin{center}\rule{0.5\linewidth}{0.5pt}\end{center}

\subsection{Instructions}\label{instructions}

\begin{enumerate}
\def\labelenumi{\arabic{enumi}.}
\tightlist
\item
  \textbf{Permitted:} lecture slides, your own notes, textbooks
  (Laplante, Liu, Barry), lab code and write-ups, a calculator.
\item
  \textbf{Not permitted:} communicating with anyone, or using the
  internet, AI assistants, or any online service.
\item
  \textbf{Show your work.} For schedulability problems write the
  formula, the substituted numbers, and the conclusion. A correct final
  answer with no working earns at most half marks. A wrong answer with
  correct method earns most of the marks.
\item
  \textbf{State any assumption} you make. Unless a question says
  otherwise:

  \begin{itemize}
  \tightlist
  \item
    scheduling is \textbf{preemptive, fixed-priority} (a larger priority
    number is \emph{not} implied --- ``higher priority'' is always
    stated in words);
  \item
    context-switch and kernel overhead are \textbf{zero};
  \item
    tasks are \textbf{independent} and \textbf{synchronous}
    (\(\Phi_i = 0\)), time unit is \textbf{ms};
  \item
    ties in priority or deadline are broken in favour of the
    \textbf{lower task index}.
  \end{itemize}
\item
  Code is FreeRTOS on a Cortex-M33 (FRDM-MCXN236). You do not need to
  remember exact header names; correct API names and argument logic are
  what matter.
\item
  Budget your time. A suggested allocation is printed beside each part.
  Do not spend more than 10 minutes stuck on one question --- move on
  and return.
\end{enumerate}

\subsubsection{Score summary}\label{score-summary}

{\def\LTcaptype{none} % do not increment counter
\begin{longtable}[]{@{}llrr@{}}
\toprule\noalign{}
Part & Topic & Points & Suggested time \\
\midrule\noalign{}
\endhead
\bottomrule\noalign{}
\endlastfoot
A & Concepts & 12 & 15 min \\
B & Task model and timing metrics & 12 & 20 min \\
C & Rate Monotonic Scheduling and RTA & 14 & 25 min \\
D & EDF and RMS vs.~EDF & 14 & 30 min \\
E & Aperiodic servers and sporadic tasks & 13 & 25 min \\
F & Queues, semaphores, mutexes & 20 & 35 min \\
G & Integrated design problem & 15 & 30 min \\
& \textbf{Total} & \textbf{100} & \textbf{180 min} \\
\end{longtable}
}

\subsubsection{Reference values
(provided)}\label{reference-values-provided}

Liu \& Layland bound \(U_{\text{lub}}(n) = n\left(2^{1/n}-1\right)\):

{\def\LTcaptype{none} % do not increment counter
\begin{longtable}[]{@{}
  >{\raggedright\arraybackslash}p{(\linewidth - 14\tabcolsep) * \real{0.1250}}
  >{\raggedright\arraybackslash}p{(\linewidth - 14\tabcolsep) * \real{0.1250}}
  >{\raggedright\arraybackslash}p{(\linewidth - 14\tabcolsep) * \real{0.1250}}
  >{\raggedright\arraybackslash}p{(\linewidth - 14\tabcolsep) * \real{0.1250}}
  >{\raggedright\arraybackslash}p{(\linewidth - 14\tabcolsep) * \real{0.1250}}
  >{\raggedright\arraybackslash}p{(\linewidth - 14\tabcolsep) * \real{0.1250}}
  >{\raggedright\arraybackslash}p{(\linewidth - 14\tabcolsep) * \real{0.1250}}
  >{\raggedright\arraybackslash}p{(\linewidth - 14\tabcolsep) * \real{0.1250}}@{}}
\toprule\noalign{}
\begin{minipage}[b]{\linewidth}\raggedright
\(n\)
\end{minipage} & \begin{minipage}[b]{\linewidth}\raggedright
1
\end{minipage} & \begin{minipage}[b]{\linewidth}\raggedright
2
\end{minipage} & \begin{minipage}[b]{\linewidth}\raggedright
3
\end{minipage} & \begin{minipage}[b]{\linewidth}\raggedright
4
\end{minipage} & \begin{minipage}[b]{\linewidth}\raggedright
5
\end{minipage} & \begin{minipage}[b]{\linewidth}\raggedright
6
\end{minipage} & \begin{minipage}[b]{\linewidth}\raggedright
\(\infty\)
\end{minipage} \\
\midrule\noalign{}
\endhead
\bottomrule\noalign{}
\endlastfoot
\(U_{\text{lub}}(n)\) & 1.0000 & 0.8284 & 0.7798 & 0.7568 & 0.7435 &
0.7348 & \(\ln 2 = 0.6931\) \\
\end{longtable}
}

Response-time recurrence:
\(R_i^{(0)} = C_i,\;\; R_i^{(k+1)} = C_i + \sum_{j \in \text{hp}(i)} \left\lceil R_i^{(k)} / T_j \right\rceil C_j\)

Demand bound:
\(\text{dbf}(L) = \sum_i \max\!\left(0,\ \left\lfloor \dfrac{L-D_i}{T_i} \right\rfloor + 1\right) C_i \le L\)

Deferrable-server bound (\(n\) periodic tasks + one DS of utilization
\(U_s\)):
\(U_{\text{lub}} = U_s + n\left[\left(\dfrac{U_s+2}{2U_s+1}\right)^{1/n} - 1\right]\)

\begin{center}\rule{0.5\linewidth}{0.5pt}\end{center}

\subsection{Part A --- Concepts (12 points · \textasciitilde15
min)}\label{part-a-concepts-12-points-15-min}

\textbf{A1. (2 pts)} Classify each system's timing requirement as
\textbf{hard}, \textbf{firm}, or \textbf{soft}, and give a one-phrase
justification for each.

\begin{enumerate}
\def\labelenumi{\alph{enumi}.}
\tightlist
\item
  Airbag deployment controller
\item
  Video-conference frame renderer (a late frame is dropped, the next one
  is shown)
\item
  Weather-station display refreshing every 5 s
\item
  Sensor-fusion sample where a stale sample is discarded and useless,
  but a miss does not endanger anyone
\end{enumerate}

\textbf{A2. (2 pts)} A colleague says: \emph{``Our MCU runs at 150 MHz
and averages 2 µs response to an interrupt, so it is a real-time
system.''} a. Explain why this statement is not sufficient. (1 pt) b.
Name \textbf{two} kernel properties that distinguish an RTOS from a
general-purpose OS. (1 pt)

\textbf{A3. (2 pts)} Define \textbf{periodic}, \textbf{sporadic}, and
\textbf{aperiodic} tasks by their arrival pattern. Which of the three
can sensibly carry a \textbf{hard} deadline \emph{and} needs no server
to be analysed, and why?

\textbf{A4. (2 pts)} Decide whether each statement is \textbf{true or
false}, with one sentence of justification.

\begin{enumerate}
\def\labelenumi{\alph{enumi}.}
\tightlist
\item
  ``RMS is optimal, so if RMS cannot schedule a task set, no scheduling
  algorithm can.'' (1 pt)
\item
  ``Under the critical-instant theorem, the worst-case response time of
  \(\tau_i\) occurs when \(\tau_i\) is released at the same instant as
  all \textbf{lower}-priority tasks.'' (1 pt)
\end{enumerate}

\textbf{A5. (2 pts)} Apply the utilization-bound test to each case below
under RMS. Answer \emph{schedulable}, \emph{inconclusive}, or
\emph{infeasible}.

{\def\LTcaptype{none} % do not increment counter
\begin{longtable}[]{@{}lll@{}}
\toprule\noalign{}
& \(n\) & \(U\) \\
\midrule\noalign{}
\endhead
\bottomrule\noalign{}
\endlastfoot
(i) & 3 & 0.75 \\
(ii) & 4 & 0.80 \\
(iii) & 2 & 1.05 \\
\end{longtable}
}

\textbf{A6. (2 pts)} State \textbf{three} differences between a FreeRTOS
\textbf{mutex} and a \textbf{binary semaphore}, and say which one you
would use to protect an I²C bus shared by three tasks.

\begin{center}\rule{0.5\linewidth}{0.5pt}\end{center}

\subsection{Part B --- Task Model and Timing Metrics (12 points ·
\textasciitilde20
min)}\label{part-b-task-model-and-timing-metrics-12-points-20-min}

\textbf{B1. (4 pts)} A job is released at \(r = 6\) ms with relative
deadline \(D = 10\) ms and execution time \(C = 4\) ms. It first gets
the CPU at \(t = 9\) and completes at \(f = 19\).

\begin{enumerate}
\def\labelenumi{\alph{enumi}.}
\tightlist
\item
  Absolute deadline and response time. (1 pt)
\item
  Lateness and tardiness. (1 pt)
\item
  Laxity at release, and laxity at \(t = 9\) (just before it starts
  running). (1 pt)
\item
  Another job of the same task, released at \(r = 6\), completes at
  \(f = 12\). Give its lateness and tardiness, and explain why lateness
  may be negative but tardiness may not. (1 pt)
\end{enumerate}

\textbf{B2. (4 pts)} A task is specified by
\(\tau = (\Phi, T, C, D) = (3,\ 10,\ 2,\ 7)\).

\begin{enumerate}
\def\labelenumi{\alph{enumi}.}
\tightlist
\item
  Give the release time and absolute deadline of its first three jobs.
  (1 pt)
\item
  Is the deadline model \emph{implicit}, \emph{constrained}, or
  \emph{arbitrary}? Which fixed-priority assignment rule applies to a
  set of such tasks? (1 pt)
\item
  Three tasks have periods \(T_a = 6,\ T_b = 10,\ T_c = 15\) ms. Compute
  the hyperperiod \(H\) and the number of jobs of each task inside one
  hyperperiod. (2 pts)
\end{enumerate}

\textbf{B3. (4 pts)} A control task must run every 10 ms. Its body takes
3 ms to execute.

\begin{Shaded}
\begin{Highlighting}[]
\DataTypeTok{void}\NormalTok{ vControlTask}\OperatorTok{(}\DataTypeTok{void} \OperatorTok{*}\NormalTok{pv}\OperatorTok{)}
\OperatorTok{\{}
    \ControlFlowTok{for} \OperatorTok{(;;)} \OperatorTok{\{}
\NormalTok{        ReadSensorAndCompute}\OperatorTok{();}            \CommentTok{/* 3 ms of work */}
\NormalTok{        vTaskDelay}\OperatorTok{(}\NormalTok{pdMS\_TO\_TICKS}\OperatorTok{(}\DecValTok{10}\OperatorTok{));}
    \OperatorTok{\}}
\OperatorTok{\}}
\end{Highlighting}
\end{Shaded}

\begin{enumerate}
\def\labelenumi{\alph{enumi}.}
\tightlist
\item
  With a 1 kHz tick, what period does this task actually have, and what
  happens over many cycles if a higher-priority task occasionally
  preempts it? (1.5 pts)
\item
  Rewrite the loop so the task is released every 10 ms without drift.
  (1.5 pts)
\item
  After your fix, can the task still start late? Explain in terms of the
  timing metrics from Week 2. (1 pt)
\end{enumerate}

\begin{center}\rule{0.5\linewidth}{0.5pt}\end{center}

\subsection{Part C --- Rate Monotonic Scheduling and RTA (14 points ·
\textasciitilde25
min)}\label{part-c-rate-monotonic-scheduling-and-rta-14-points-25-min}

Task set \textbf{P} (priorities assigned by RMS):

{\def\LTcaptype{none} % do not increment counter
\begin{longtable}[]{@{}lrr@{}}
\toprule\noalign{}
Task & \(C_i\) & \(T_i = D_i\) \\
\midrule\noalign{}
\endhead
\bottomrule\noalign{}
\endlastfoot
\(\tau_1\) & 1 & 4 \\
\(\tau_2\) & 2 & 6 \\
\(\tau_3\) & 3 & 12 \\
\end{longtable}
}

\textbf{C1. (3 pts)} Assign RMS priorities. Compute \(U\) and
\(U_{\text{lub}}(3)\) and state what the utilization test concludes.

\textbf{C2. (6 pts)} Perform response-time analysis for all three tasks.
\textbf{Show every iteration} of the recurrence. State whether set P is
schedulable.

\textbf{C3. (3 pts)} Draw the RMS Gantt chart for \(t \in [0, 12]\).
Mark every release, every deadline, and every preemption. Your chart
must agree with the response times from C2.

\textbf{C4. (2 pts)} Keeping \(\tau_1\) and \(\tau_2\) unchanged, find
the \textbf{largest integer value of \(C_3\)} for which the set is still
schedulable under RMS. Justify with RTA. Is the Liu \& Layland bound a
necessary condition? Explain in one sentence.

\begin{center}\rule{0.5\linewidth}{0.5pt}\end{center}

\subsection{Part D --- EDF and RMS vs.~EDF (14 points ·
\textasciitilde30
min)}\label{part-d-edf-and-rms-vs.-edf-14-points-30-min}

Task set \textbf{Q} (implicit deadlines):

{\def\LTcaptype{none} % do not increment counter
\begin{longtable}[]{@{}lrr@{}}
\toprule\noalign{}
Task & \(C_i\) & \(T_i = D_i\) \\
\midrule\noalign{}
\endhead
\bottomrule\noalign{}
\endlastfoot
\(\tau_1\) & 2 & 6 \\
\(\tau_2\) & 2 & 8 \\
\(\tau_3\) & 5 & 12 \\
\end{longtable}
}

\textbf{D1. (2 pts)} Compute \(U\). What does the EDF theorem say about
set Q? What does the Liu \& Layland test say about it under RMS?

\textbf{D2. (3 pts)} Run RTA on set Q under RMS. Show the iterations and
state the outcome.

\textbf{D3. (6 pts)} Construct the \textbf{EDF schedule} for set Q over
one hyperperiod. Present it as a table of execution intervals (for
example \texttt{{[}0,2{]}\ J1,1}), where \(J_{i,k}\) is the \(k\)-th job
of \(\tau_i\). Then list the completion time of \textbf{every} job and
confirm each meets its deadline. Use the tie-break rule from the
instructions.

\textbf{D4. (3 pts)} Constrained-deadline set R:

{\def\LTcaptype{none} % do not increment counter
\begin{longtable}[]{@{}lrrr@{}}
\toprule\noalign{}
Task & \(C_i\) & \(T_i\) & \(D_i\) \\
\midrule\noalign{}
\endhead
\bottomrule\noalign{}
\endlastfoot
\(\tau_1\) & 3 & 8 & 4 \\
\(\tau_2\) & 3 & 10 & 5 \\
\end{longtable}
}

\begin{enumerate}
\def\labelenumi{\alph{enumi}.}
\tightlist
\item
  Compute \(U\). Does \(U \le 1\) guarantee EDF schedulability here? (1
  pt)
\item
  Compute \(\text{dbf}(4)\) and \(\text{dbf}(5)\). (1 pt)
\item
  Is set R EDF-schedulable? Explain what goes wrong in physical terms
  (what must happen before \(t = 5\)?). (1 pt)
\end{enumerate}

\begin{center}\rule{0.5\linewidth}{0.5pt}\end{center}

\subsection{Part E --- Aperiodic Servers and Sporadic Tasks (13 points ·
\textasciitilde25
min)}\label{part-e-aperiodic-servers-and-sporadic-tasks-13-points-25-min}

\textbf{E1. (5 pts)} Periodic tasks \(\tau_1 = (C{=}1, T{=}5)\) and
\(\tau_2 = (C{=}2, T{=}10)\) share the CPU with an aperiodic server of
capacity \(C_s = 1\) and period \(T_s = 4\). Priorities are by RMS, so
the server has the \textbf{highest} priority.

\begin{enumerate}
\def\labelenumi{\alph{enumi}.}
\tightlist
\item
  If the server is a \textbf{polling server}, apply the Liu \& Layland
  test with the server counted as a periodic task. What is the verdict?
  (2 pts)
\item
  If the server is a \textbf{deferrable server}, compute the DS
  utilization bound for \(n = 2\) periodic tasks and \(U_s = C_s/T_s\).
  Is the set still guaranteed schedulable? Why can the ordinary
  \(U_{\text{lub}}(3)\) not be used? (3 pts)
\end{enumerate}

\textbf{E2. (3 pts)} In the system of E1, an aperiodic request needing 1
ms of service arrives at \(t = 1\). The server's poll at \(t = 0\) found
the queue empty.

\begin{enumerate}
\def\labelenumi{\alph{enumi}.}
\tightlist
\item
  When does the request complete, and what is its response time, under
  the \textbf{polling server}? (1 pt)
\item
  Same question under the \textbf{deferrable server}. (1 pt)
\item
  The \textbf{sporadic server} gives the same response time as (b). Why
  can it nevertheless be analysed as an ordinary periodic task while the
  deferrable server cannot? (1 pt)
\end{enumerate}

\textbf{E3. (5 pts)} A system has three activities:

{\def\LTcaptype{none} % do not increment counter
\begin{longtable}[]{@{}llrrr@{}}
\toprule\noalign{}
Activity & Behaviour & \(C\) & Period / min. inter-arrival & \(D\) \\
\midrule\noalign{}
\endhead
\bottomrule\noalign{}
\endlastfoot
Control loop \(\tau_1\) & periodic & 2 & \(T = 10\) & 10 \\
Fault alarm \(\tau_2\) & \textbf{sporadic} & 2 & \(a = 20\) &
\textbf{3} \\
Logger \(\tau_3\) & periodic & 3 & \(T = 20\) & 20 \\
\end{longtable}
}

\begin{enumerate}
\def\labelenumi{\alph{enumi}.}
\tightlist
\item
  How do you model the sporadic alarm for analysis, and why is this safe
  for any slower arrival pattern? (1 pt)
\item
  Which priority-assignment rule must you use, and what is the resulting
  priority order? (1 pt)
\item
  Run RTA for all three tasks with that assignment. Is the set
  schedulable? (2 pts)
\item
  Show that plain \textbf{RMS} would have failed for the alarm. (1 pt)
\end{enumerate}

\begin{center}\rule{0.5\linewidth}{0.5pt}\end{center}

\subsection{Part F --- Queues, Semaphores, Mutexes (20 points ·
\textasciitilde35
min)}\label{part-f-queues-semaphores-mutexes-20-points-35-min}

\textbf{F1. (5 pts)} The code below logs bytes received by a UART ISR;
the logger and a heartbeat task both print on the same UART and are
meant to be serialised by a mutex. It contains \textbf{at least five
distinct defects} (functional bugs or incorrect API use). Find
\textbf{five}, explain the failure each causes, and give the fix. (1 pt
each.)

\begin{Shaded}
\begin{Highlighting}[]
\DataTypeTok{static}\NormalTok{ SemaphoreHandle\_t xUartMutex}\OperatorTok{;}
\DataTypeTok{static}\NormalTok{ QueueHandle\_t     xRxQueue}\OperatorTok{;}

\DataTypeTok{void}\NormalTok{ UART\_IRQHandler}\OperatorTok{(}\DataTypeTok{void}\OperatorTok{)}
\OperatorTok{\{}
    \DataTypeTok{uint8\_t}\NormalTok{ b }\OperatorTok{=}\NormalTok{ UART\_ReadByte}\OperatorTok{();}
\NormalTok{    xQueueSendToBack}\OperatorTok{(}\NormalTok{xRxQueue}\OperatorTok{,} \OperatorTok{\&}\NormalTok{b}\OperatorTok{,}\NormalTok{ portMAX\_DELAY}\OperatorTok{);}
\OperatorTok{\}}

\DataTypeTok{void}\NormalTok{ vLoggerTask}\OperatorTok{(}\DataTypeTok{void} \OperatorTok{*}\NormalTok{pv}\OperatorTok{)}
\OperatorTok{\{}
    \DataTypeTok{uint8\_t}\NormalTok{ b}\OperatorTok{;}
    \ControlFlowTok{for} \OperatorTok{(;;)} \OperatorTok{\{}
\NormalTok{        xQueueReceive}\OperatorTok{(}\NormalTok{xRxQueue}\OperatorTok{,} \OperatorTok{\&}\NormalTok{b}\OperatorTok{,}\NormalTok{ portMAX\_DELAY}\OperatorTok{);}
\NormalTok{        xSemaphoreTake}\OperatorTok{(}\NormalTok{xUartMutex}\OperatorTok{,}\NormalTok{ portMAX\_DELAY}\OperatorTok{);}
\NormalTok{        PRINTF}\OperatorTok{(}\StringTok{"rx=0x}\SpecialCharTok{\%02x\textbackslash{}r\textbackslash{}n}\StringTok{"}\OperatorTok{,}\NormalTok{ b}\OperatorTok{);}
        \ControlFlowTok{if} \OperatorTok{(}\NormalTok{b }\OperatorTok{==} \BaseNTok{0x1B}\OperatorTok{)} \OperatorTok{\{}                                     \CommentTok{/* ESC */}
\NormalTok{            PRINTF}\OperatorTok{(}\StringTok{"escape received}\SpecialCharTok{\textbackslash{}r\textbackslash{}n}\StringTok{"}\OperatorTok{);}
            \ControlFlowTok{continue}\OperatorTok{;}
        \OperatorTok{\}}
\NormalTok{        xSemaphoreGive}\OperatorTok{(}\NormalTok{xUartMutex}\OperatorTok{);}
    \OperatorTok{\}}
\OperatorTok{\}}

\DataTypeTok{void}\NormalTok{ vHeartbeatTask}\OperatorTok{(}\DataTypeTok{void} \OperatorTok{*}\NormalTok{pv}\OperatorTok{)}
\OperatorTok{\{}
    \ControlFlowTok{for} \OperatorTok{(;;)} \OperatorTok{\{}
\NormalTok{        PRINTF}\OperatorTok{(}\StringTok{"alive}\SpecialCharTok{\textbackslash{}r\textbackslash{}n}\StringTok{"}\OperatorTok{);}
\NormalTok{        vTaskDelay}\OperatorTok{(}\NormalTok{pdMS\_TO\_TICKS}\OperatorTok{(}\DecValTok{1000}\OperatorTok{));}
    \OperatorTok{\}}
\OperatorTok{\}}

\DataTypeTok{int}\NormalTok{ main}\OperatorTok{(}\DataTypeTok{void}\OperatorTok{)}
\OperatorTok{\{}
\NormalTok{    BOARD\_InitHardware}\OperatorTok{();}
\NormalTok{    xRxQueue   }\OperatorTok{=}\NormalTok{ xQueueCreate}\OperatorTok{(}\DecValTok{8}\OperatorTok{,} \KeywordTok{sizeof}\OperatorTok{(}\DataTypeTok{uint32\_t}\OperatorTok{));}
\NormalTok{    xUartMutex }\OperatorTok{=}\NormalTok{ xSemaphoreCreateBinary}\OperatorTok{();}
\NormalTok{    xTaskCreate}\OperatorTok{(}\NormalTok{vLoggerTask}\OperatorTok{,}    \StringTok{"log"}\OperatorTok{,} \DecValTok{256}\OperatorTok{,}\NormalTok{ NULL}\OperatorTok{,} \DecValTok{1}\OperatorTok{,}\NormalTok{ NULL}\OperatorTok{);}
\NormalTok{    xTaskCreate}\OperatorTok{(}\NormalTok{vHeartbeatTask}\OperatorTok{,} \StringTok{"hb"}\OperatorTok{,}  \DecValTok{256}\OperatorTok{,}\NormalTok{ NULL}\OperatorTok{,} \DecValTok{1}\OperatorTok{,}\NormalTok{ NULL}\OperatorTok{);}
\NormalTok{    vTaskStartScheduler}\OperatorTok{();}
    \ControlFlowTok{for} \OperatorTok{(;;)} \OperatorTok{\{\}}
\OperatorTok{\}}
\end{Highlighting}
\end{Shaded}

\textbf{F2. (5 pts)} For each scenario choose the \textbf{most
appropriate primitive} from: \emph{queue}, \emph{binary semaphore},
\emph{counting semaphore}, \emph{mutex}, \emph{critical section}. Give a
one-sentence reason and name the key API call(s). (1 pt each.)

\begin{enumerate}
\def\labelenumi{\alph{enumi}.}
\tightlist
\item
  A DMA-complete ISR must wake a handler task. No data is passed.
\item
  Three tasks take turns using one I²C bus; each transaction lasts about
  200 µs of task-level code.
\item
  A pool of four identical DMA channels is shared by many tasks.
\item
  A sensor task sends a 12-byte struct \texttt{\{timestamp,\ x,\ y\}} to
  a filter task. Bursts of up to three samples can occur while the
  filter is busy.
\item
  A timer ISR and a task both execute \texttt{gTicks++} on a 32-bit
  variable. The task's access is two instructions long.
\end{enumerate}

\textbf{F3. (5 pts)} A queue \texttt{q} of length \textbf{2} carries
\texttt{int} items. Task \textbf{P} has priority \textbf{2} and task
\textbf{C} has priority \textbf{1} (higher number = higher priority for
this question only). Both are created before the scheduler starts.
Assume all calls other than blocking and \texttt{vTaskDelay} take zero
time, and \texttt{PRINTF} is instantaneous.

\begin{Shaded}
\begin{Highlighting}[]
\DataTypeTok{void}\NormalTok{ vP}\OperatorTok{(}\DataTypeTok{void} \OperatorTok{*}\NormalTok{pv}\OperatorTok{)} \OperatorTok{\{}                  \DataTypeTok{void}\NormalTok{ vC}\OperatorTok{(}\DataTypeTok{void} \OperatorTok{*}\NormalTok{pv}\OperatorTok{)} \OperatorTok{\{}
    \ControlFlowTok{for} \OperatorTok{(}\DataTypeTok{int}\NormalTok{ i }\OperatorTok{=} \DecValTok{1}\OperatorTok{;}\NormalTok{ i }\OperatorTok{\textless{}=} \DecValTok{4}\OperatorTok{;}\NormalTok{ i}\OperatorTok{++)}         \ControlFlowTok{for} \OperatorTok{(;;)} \OperatorTok{\{}
\NormalTok{        xQueueSendToBack}\OperatorTok{(}\NormalTok{q}\OperatorTok{,} \OperatorTok{\&}\NormalTok{i}\OperatorTok{,}              \DataTypeTok{int}\NormalTok{ v}\OperatorTok{;}
\NormalTok{                         portMAX\_DELAY}\OperatorTok{);}\NormalTok{     xQueueReceive}\OperatorTok{(}\NormalTok{q}\OperatorTok{,} \OperatorTok{\&}\NormalTok{v}\OperatorTok{,}\NormalTok{ portMAX\_DELAY}\OperatorTok{);}
\NormalTok{    vTaskDelete}\OperatorTok{(}\NormalTok{NULL}\OperatorTok{);}\NormalTok{                       PRINTF}\OperatorTok{(}\StringTok{"}\SpecialCharTok{\%d}\StringTok{ @}\SpecialCharTok{\%u\textbackslash{}n}\StringTok{"}\OperatorTok{,}\NormalTok{ v}\OperatorTok{,}\NormalTok{ xTaskGetTickCount}\OperatorTok{());}
\OperatorTok{\}}\NormalTok{                                            vTaskDelay}\OperatorTok{(}\DecValTok{10}\OperatorTok{);}
                                         \OperatorTok{\}}
                                     \OperatorTok{\}}
\end{Highlighting}
\end{Shaded}

\begin{enumerate}
\def\labelenumi{\alph{enumi}.}
\tightlist
\item
  Give the exact output (value and tick) printed by C. (2 pts)
\item
  What are the contents of the queue at tick 5? (1 pt)
\item
  State at which instants task P blocks and unblocks, and explain why C
  is preempted inside its first \texttt{xQueueReceive} call. (2 pts)
\end{enumerate}

\textbf{F4. (5 pts)}

\begin{enumerate}
\def\labelenumi{\alph{enumi}.}
\item
  Two tasks each execute \texttt{gCounter++} once on a shared
  \texttt{volatile\ uint32\_t} with initial value 10. On a load-store
  architecture the increment is
  \texttt{LDR\ r0,{[}gCounter{]};\ ADD\ r0,r0,\#1;\ STR\ r0,{[}gCounter{]}}.
  Write an interleaving of the six instructions that leaves
  \texttt{gCounter\ ==\ 11}. Say where the preemption happened. (2 pts)
\item
  A function \texttt{vLog()} takes \texttt{xMutex} (created with
  \texttt{xSemaphoreCreateMutex()}) and, while holding it, calls
  \texttt{vFormat()}, which also takes \texttt{xMutex}. What happens
  when a task runs this code with \texttt{portMAX\_DELAY}? Give
  \textbf{two} distinct ways to fix it. (2 pts)
\item
  Why can a mutex not be given from an ISR, and why is a binary
  semaphore (not a mutex) the right object for ISR → task signalling? (1
  pt)
\end{enumerate}

\begin{center}\rule{0.5\linewidth}{0.5pt}\end{center}

\subsection{Part G --- Integrated Design Problem (15 points ·
\textasciitilde30
min)}\label{part-g-integrated-design-problem-15-points-30-min}

You are designing the firmware for a small motor-drive controller.
Ignore context-switch cost and ignore blocking caused by the shared UART
mutex (those effects are analysed in the next module).

{\def\LTcaptype{none} % do not increment counter
\begin{longtable}[]{@{}
  >{\raggedright\arraybackslash}p{(\linewidth - 8\tabcolsep) * \real{0.1667}}
  >{\raggedright\arraybackslash}p{(\linewidth - 8\tabcolsep) * \real{0.1667}}
  >{\raggedleft\arraybackslash}p{(\linewidth - 8\tabcolsep) * \real{0.2222}}
  >{\raggedleft\arraybackslash}p{(\linewidth - 8\tabcolsep) * \real{0.2222}}
  >{\raggedleft\arraybackslash}p{(\linewidth - 8\tabcolsep) * \real{0.2222}}@{}}
\toprule\noalign{}
\begin{minipage}[b]{\linewidth}\raggedright
Activity
\end{minipage} & \begin{minipage}[b]{\linewidth}\raggedright
Description
\end{minipage} & \begin{minipage}[b]{\linewidth}\raggedleft
\(C\) (ms)
\end{minipage} & \begin{minipage}[b]{\linewidth}\raggedleft
Period / min. gap
\end{minipage} & \begin{minipage}[b]{\linewidth}\raggedleft
\(D\) (ms)
\end{minipage} \\
\midrule\noalign{}
\endhead
\bottomrule\noalign{}
\endlastfoot
\textbf{E-stop handler} & Runs when a button ISR fires; at most once per
50 ms & 1 & \(a = 50\) & 5 \\
\textbf{Current loop} & Reads ADC, updates PWM & 1 & \(T = 5\) & 5 \\
\textbf{CLI} & Operator commands from the serial console, arbitrary
arrival, soft & --- & --- & --- \\
\textbf{Speed loop} & Computes speed, sends a result struct to the
logger & 2 & \(T = 20\) & 20 \\
\textbf{Logger} & Receives results, prints them over the UART & 4 &
\(T = 40\) & 40 \\
\end{longtable}
}

CLI requests are served by a \textbf{polling server} with \(C_s = 2\)
ms, \(T_s = 10\) ms. The CLI handler and the logger both print on the
same UART.

\textbf{G1. (2 pts)} For each of the five activities say whether it is
periodic, sporadic, or aperiodic, and whether its deadline is hard,
firm, or soft.

\textbf{G2. (3 pts)} Give the full priority order (highest to lowest) of
the five schedulable entities (including the server). Name the rule you
use, and justify the placement of the E-stop handler relative to the
current loop (they have equal deadlines).

\textbf{G3. (2 pts)} Compute total utilization (including the server)
and apply the utilization-bound test. What does it conclude?

\textbf{G4. (5 pts)} Compute the worst-case response times of the
\textbf{speed loop} and of the \textbf{logger} with RTA (show the
iterations). Compare each with its deadline.

\textbf{G5. (3 pts)} Identify the three inter-task communication objects
this design needs. For each, state the primitive, who produces and who
consumes, and which API variant (task or ISR) is called from where:

\begin{enumerate}
\def\labelenumi{\alph{enumi}.}
\tightlist
\item
  button ISR → E-stop handler
\item
  speed loop → logger
\item
  logger and CLI handler → UART
\end{enumerate}

\begin{center}\rule{0.5\linewidth}{0.5pt}\end{center}

\emph{End of examination.}

\end{document}
